%% ma: L11-B10-0003
%% lop: 11
%% chuong: 4
%% bai: 10
%% dang: 11.10.3
%% loai: TN4
%% muc: TH
%% dap_an: "D"
%% nguon: "de-cuong-giua-ky-1-toan-11-nam-2026-2027-truong-thpt-hoang-van-thu-ha-noi-2-8.pdf (M08, Đề cương giữa kì 1 Toán 11 THPT Hoàng Văn Thụ, Hà Nội), Phần 1 Câu 25"
%% kien_thuc: "Bài 10 (giao điểm của đường thẳng và mặt phẳng)"
%% trang_thai: chinh-thuc
%% ghi_chu: "Đã duyệt 29/10/2026."
\begin{ex}
Cho tứ diện $ABCD$. Gọi $E$, $F$ lần lượt là trung điểm của các cạnh $AB$, $CD$. Điểm $G$ là trọng tâm tam giác $BCD$. Khi đó giao điểm của đường thẳng $EG$ và mặt phẳng $(ACD)$ là
\choice{điểm $F$}{giao điểm của đường thẳng $EG$ và $CD$}{giao điểm của đường thẳng $EG$ và $AC$}{\True giao điểm của đường thẳng $EG$ và $AF$}
\loigiai{$G$ thuộc trung tuyến $BF$ nên $EG\subset(ABF)$. Mà $(ABF)\cap(ACD)=AF$, nên giao điểm của $EG$ và $(ACD)$ là giao điểm của $EG$ và $AF$.}
\end{ex}
